\documentclass[aspectratio=169]{beamer}

% ---------- ENCODING & FONTS ----------
\usepackage[utf8]{inputenc}
\usepackage[T1]{fontenc}

% ---------- THEME & LOOK ----------
\usetheme{Madrid}
\setbeamertemplate{navigation symbols}{}
\setbeamercovered{transparent}
\setbeamertemplate{itemize items}[circle]
\setbeamertemplate{enumerate items}[default]
\setlength{\parskip}{0.25em}

% ---------- COLORS ----------
\definecolor{UNBCGreen}{RGB}{0,102,51}
\definecolor{UNBCDark}{RGB}{0,74,37}
\definecolor{UNBCGold}{RGB}{189,155,96}
\definecolor{SoftCream}{RGB}{249,247,242}
\definecolor{SoftGray}{RGB}{244,246,248}
\definecolor{AccentBlue}{RGB}{41,98,255}
\definecolor{AccentOrange}{RGB}{230,126,34}
\definecolor{AccentTeal}{RGB}{0,150,136}
\definecolor{AccentRed}{RGB}{192,57,43}
% legacy aliases kept for in-slide content
\definecolor{accent}{RGB}{0,102,51}
\definecolor{soft}{RGB}{244,246,248}
\definecolor{canadared}{RGB}{196,30,58}
\definecolor{leafgreen}{RGB}{0,130,80}
\definecolor{goldenrod}{RGB}{189,155,96}
\definecolor{deeporange}{RGB}{220,90,20}
\setbeamercolor{normal text}{fg=black,bg=white}
\setbeamercolor{structure}{fg=UNBCDark}
\setbeamercolor{title}{fg=white,bg=UNBCDark}
\setbeamercolor{subtitle}{fg=UNBCGold}
\setbeamercolor{frametitle}{fg=white,bg=UNBCDark}
\setbeamercolor{block title}{fg=white,bg=UNBCDark}
\setbeamercolor{block body}{fg=black,bg=SoftGray}
\setbeamercolor{block title alerted}{fg=white,bg=AccentRed}
\setbeamercolor{block body alerted}{fg=black,bg=SoftGray}
\setbeamercolor{item}{fg=UNBCDark}
\setbeamercolor{section in toc}{fg=UNBCDark}
\setbeamercolor{alerted text}{fg=AccentRed}
\setbeamercolor{author in head/foot}{fg=white,bg=UNBCDark}
\setbeamercolor{date in head/foot}{fg=white,bg=UNBCDark}

% ---------- PACKAGES ----------
\usepackage{booktabs}
\usepackage{amsmath, amssymb}
\usepackage{array}
\usepackage{hyperref}
\usepackage{tikz}
\hypersetup{colorlinks=true,linkcolor=UNBCDark,urlcolor=UNBCDark}

% ---------- SAFE TRANSITION FALLBACK ----------
\makeatletter
\@ifundefined{transfade}{\newcommand{\transfade}[1][]{}}{}
\makeatother

% ---------- UNBC BRANDING ----------
\usepackage{graphicx}
\usepackage{lmodern}
\usepackage{tcolorbox}
\setbeamertemplate{headline}{}
\setbeamertemplate{footline}{%
  \leavevmode%
  \hbox{%
  \begin{beamercolorbox}[wd=.78\paperwidth,ht=2.8ex,dp=1.2ex,leftskip=1em]{author in head/foot}%
  \color{white}\textbf{ECON 101} \hspace{0.5em} \textcolor{UNBCGold}{|} \hspace{0.5em} Macroeconomics%
  \end{beamercolorbox}%
  \begin{beamercolorbox}[wd=.22\paperwidth,ht=2.8ex,dp=1.2ex,rightskip=1em plus1fil]{date in head/foot}%
  \color{white}\insertframenumber/\inserttotalframenumber%
  \end{beamercolorbox}}%
}

\newtcolorbox{mybox}[2][]{
  colback=#2!8,
  colframe=#2,
  boxrule=0.8pt,
  arc=2mm,
  left=2mm,right=2mm,top=1.5mm,bottom=1.5mm,
  #1
}

\newcommand{\topbanner}[1]{%
\begin{tikzpicture}[remember picture,overlay]
\fill[UNBCGold] ([yshift=-0.65cm]current page.north west) rectangle ([yshift=-0.48cm]current page.north east);
\node[anchor=north west,font=\small\bfseries,text=UNBCDark] at ([xshift=0.35cm,yshift=-0.78cm]current page.north west) {#1};
\end{tikzpicture}%
}


% ---------- TITLE ----------
\title{\textbf{Welcome to ECON 101}}
\subtitle{{\large Macroeconomics}}
\author{Muhebullah Karimzada}
\institute{School of Economics \\ University of Northern British Columbia}
\date{September 9, 2026}

% ============================================================
\begin{document}
% ============================================================

\begin{frame}[plain]
  \titlepage
\end{frame}

% ============================================================
\begin{frame}{Today's Agenda}
\transfade
\begin{itemize}
  \item Course overview
  \item Instructor \& contact information
  \item Assessments and grading
  \item Course policies
  \item RISE modules
  \item Student resources
  \item \textbf{Introduction: What is Economics?}
  \item \textbf{Microeconomics vs.\ Macroeconomics}
  \item \textbf{Key macroeconomic questions}
\end{itemize}
\end{frame}

% ============================================================
\section{Course Overview}
% ============================================================

\begin{frame}{About Your Instructor}
\transfade
\begin{columns}
  \begin{column}{0.55\textwidth}
    \begin{block}{Contact Information}
      \begin{itemize}
        \item \textbf{Instructor:} Muhebullah Karimzada
        \item \textbf{Email:} karimzada@unbc.ca
        \item \textbf{Office:} Room 10-4522
      \end{itemize}
    \end{block}
    \vspace{0.4em}
    \begin{block}{Office Hours}
      \begin{itemize}
        \item \textbf{Tuesday:} 1:00 pm -- 2:30 pm
        \item \textbf{Wednesday:} 11:30 am -- 1:00 pm
        \item \textbf{Thursday:} 1:00 pm -- 2:00 pm
        \item Or by appointment
      \end{itemize}
    \end{block}
  \end{column}
  \begin{column}{0.42\textwidth}
    \begin{alertblock}{Teaching Assistant}
      \centering TBA
    \end{alertblock}
    \vspace{0.4em}
    \begin{block}{Lectures}
      \begin{itemize}
        \item Mon \& Wed
        \item 10:00 am -- 11:20 am
        \item Room 7-212
      \end{itemize}
    \end{block}
  \end{column}
\end{columns}
\end{frame}

% ============================================================
\begin{frame}{What is This Course About?}
\transfade
\begin{block}{Course Description}
  This course offers an in-depth understanding of the key concepts, theories, and principles that drive the functioning of modern economies at the \textbf{macroeconomic level}.
\end{block}
\vspace{0.5em}
\textbf{Topics explored include:}
\begin{columns}
  \begin{column}{0.48\textwidth}
    \begin{itemize}
      \item Market forces of supply and demand
      \item National income determination
      \item Economic growth
      \item International trade
      \item Savings and investment
    \end{itemize}
  \end{column}
  \begin{column}{0.48\textwidth}
    \begin{itemize}
      \item Unemployment and inflation
      \item Monetary and banking systems
      \item Fiscal and monetary policies
      \item Short-term economic fluctuations
    \end{itemize}
  \end{column}
\end{columns}
\end{frame}

% ============================================================
\begin{frame}{Required Textbook \& Resources}
\transfade
\begin{block}{Required Text}
  \textbf{Macroeconomics}, 4th Canadian Edition \\
  \textit{Hubbard et al.}
\end{block}
\vspace{0.5em}
\begin{block}{Additional Resources}
  \begin{itemize}
    \item Lecture slides, practice problem sets posted on \textbf{Moodle}
    \item Students are responsible for downloading and saving all course material from Moodle before access is shut off at semester end
  \end{itemize}
\end{block}
\vspace{0.3em}
\begin{alertblock}{Important}
  Keep up with the readings \textit{before} each lecture for best results.
\end{alertblock}
\end{frame}

% ============================================================
\begin{frame}{Course Schedule at a Glance}
\transfade
\footnotesize
\begin{tabular}{@{}cll@{}}
  \toprule
  \textbf{Week} & \textbf{Dates} & \textbf{Topic} \\
  \midrule
  1  & Sep 9             & \textbf{Introduction to Economics} \\
  2  & Sep 14 \& 16      & Ch.1 Foundations and Models; Ch.2 Trade-offs \& Comparative Advantage \\
  3  & Sep 21 \& 23      & Ch.3 Supply and Demand \\
  4  & Sep 28 \& 30      & Ch.4 GDP: Measuring Total Production and Income \\
  5  & Oct 5 \& 7        & Ch.5 Unemployment and Inflation \\
  6  & Oct 14            & Ch.6 Economic Growth, Financial System, Business Cycles \\
  7  & Oct 19 \& 21      & Ch.7 Long-run Economic Growth \\
  8  & Oct 26 \& 28      & Ch.8 Aggregate Expenditure and Output \\
  9  & Nov 2 \& 4        & Ch.9 Aggregate Demand and Supply Analysis \\
  --  & Nov 9--13        & \textit{UNBC Mid-Semester Break} \\
  10 & Nov 16 \& 18      & Ch.10 Money, Banks, and the Bank of Canada \\
  11 & Nov 23 \& 25      & Ch.11 Monetary Policy \\
  12 & Nov 30 \& Dec 2   & Ch.12 Fiscal Policy \\
  13 & Dec 6 \& 8        & Final Exam Review \\
  \bottomrule
\end{tabular}
\end{frame}

% ============================================================
\section{Assessments \& Grading}
% ============================================================

\begin{frame}{Assessments \& Grading}
\transfade
\begin{columns}
  \begin{column}{0.52\textwidth}
    \begin{block}{Grade Breakdown}
      \begin{tabular}{@{}lcc@{}}
        \toprule
        \textbf{Assessment} & \textbf{Weight} & \textbf{Date} \\
        \midrule
        In-class Quizzes (4 of 5) & 10\%  & Various \\
        Midterm Exam 1            & 25\%  & Oct 19 \\
        Midterm Exam 2            & 25\%  & Nov 18 \\
        Final Exam                & 40\%  & TBA \\
        \midrule
        \textbf{Total}            & \textbf{100\%} & \\
        \bottomrule
      \end{tabular}
    \end{block}
  \end{column}
  \begin{column}{0.45\textwidth}
    \begin{block}{Quiz Dates}
      \begin{itemize}
        \item Quiz 1: Sep 28 (Ch.2--3)
        \item Quiz 2: Oct 7 (Ch.4)
        \item Quiz 3: Oct 14 (Ch.5)
        \item Quiz 4: Oct 28 (Ch.8)
        \item Quiz 5: Dec 2 (Ch.10--11)
      \end{itemize}
    \end{block}
    \begin{alertblock}{Bonus}
      RISE Bonus Modules: up to \textbf{+5\%}
    \end{alertblock}
  \end{column}
\end{columns}
\end{frame}

% ============================================================
\begin{frame}{Grading Scale}
\transfade
\begin{columns}
  \begin{column}{0.48\textwidth}
    \begin{tabular}{@{}ccc@{}}
      \toprule
      \textbf{Letter} & \textbf{GPA} & \textbf{Percentage} \\
      \midrule
      A+  & 4.33 & 90--100\% \\
      A   & 4.00 & 85--89.9\% \\
      A-  & 3.67 & 80--84.9\% \\
      B+  & 3.33 & 77--79.9\% \\
      B   & 3.00 & 73--76.9\% \\
      B-  & 2.67 & 70--72.9\% \\
      \bottomrule
    \end{tabular}
  \end{column}
  \begin{column}{0.48\textwidth}
    \begin{tabular}{@{}ccc@{}}
      \toprule
      \textbf{Letter} & \textbf{GPA} & \textbf{Percentage} \\
      \midrule
      C+  & 2.33 & 67--69.9\% \\
      C   & 2.00 & 63--66.9\% \\
      C-  & 1.67 & 60--62.9\% \\
      D+  & 1.33 & 57--59.9\% \\
      D   & 1.00 & 53--56.9\% \\
      F   & 0.00 & 0--49.9\%  \\
      \bottomrule
    \end{tabular}
  \end{column}
\end{columns}
\end{frame}

% ============================================================
\section{Course Policies}
% ============================================================

\begin{frame}{Important Course Policies}
\transfade
\begin{columns}
  \begin{column}{0.48\textwidth}
    \begin{block}{Quizzes}
      \begin{itemize}
        \item Held in last 30 min of lecture
        \item Top \textbf{4 of 5} scores count
        \item One free pass --- missed quiz counts as the free pass
        \item \textbf{No make-up quizzes}
      \end{itemize}
    \end{block}
    \vspace{0.3em}
    \begin{block}{Midterms}
      \begin{itemize}
        \item \textbf{No make-up midterms}
        \item Missed midterm (with documentation): 50\% weight $\to$ other midterm; 50\% $\to$ final
        \item Must write \textbf{at least one} midterm
      \end{itemize}
    \end{block}
  \end{column}
  \begin{column}{0.48\textwidth}
    \begin{block}{All Assessments}
      \begin{itemize}
        \item Strictly \textbf{individual}
        \item \textbf{Closed book}
        \item Mix of: short answer, analytical, numerical, true/false, MCQ, fill-in-the-blank
      \end{itemize}
    \end{block}
    \vspace{0.3em}
    \begin{alertblock}{Generative AI Policy}
      Use of generative AI is \textbf{NOT allowed} in this course unless specifically authorized by the instructor.
    \end{alertblock}
  \end{column}
\end{columns}
\end{frame}

% ============================================================
\begin{frame}{RISE Modules}
\transfade
\begin{block}{What is FY RISE?}
  The First-Year RISE Program helps you develop practical study and learning skills for success at UNBC.
\end{block}
\vspace{0.3em}
\begin{columns}
  \begin{column}{0.48\textwidth}
    \begin{block}{Required Modules \textnormal{(due Sep 25, 2026)}}
      \begin{itemize}
        \item Time Management
        \item Academic Reading
        \item Effective Studying
      \end{itemize}
      \textbf{Required to write quizzes and exams.}
    \end{block}
  \end{column}
  \begin{column}{0.48\textwidth}
    \begin{block}{Bonus Modules \textnormal{(due Nov 1, 2026)}}
      \begin{itemize}
        \item Academic Habits of Mind
        \item Generative AI in University
      \end{itemize}
      \textbf{+2.5\% per module (max +5\%)}
    \end{block}
  \end{column}
\end{columns}
\vspace{0.3em}
\begin{alertblock}{Already completed a required module?}
  Email your RISE badge number to \texttt{karimzada@unbc.ca} --- no need to redo it.
\end{alertblock}
\end{frame}

% ============================================================
\begin{frame}{Student Resources at UNBC}
\transfade
\begin{columns}
  \begin{column}{0.48\textwidth}
    \begin{block}{Academic Support}
      \begin{itemize}
        \item \textbf{Academic Success Centre (ASC)} --- free tutoring, study skills \\ Room 5-139 $|$ asc@unbc.ca
        \item \textbf{Academic Advising} --- course load, degree requirements \\ Room 7-702 $|$ advising@unbc.ca
        \item \textbf{Access Resource Centre (ARC)} --- academic accommodations \\ arc@unbc.ca
      \end{itemize}
    \end{block}
  \end{column}
  \begin{column}{0.48\textwidth}
    \begin{block}{Wellbeing \& Safety}
      \begin{itemize}
        \item \textbf{UNBC Counselling Services} \\ wellness@unbc.ca
        \item \textbf{Crisis Centre Northern BC (24/7)} \\ 1-888-562-1214
        \item \textbf{UNBC SAFE App} --- emergency contacts \& safety tips
        \item \textbf{NUGSS Food Bank} --- Building 6, open to all students
      \end{itemize}
    \end{block}
  \end{column}
\end{columns}
\end{frame}

% ============================================================
\section{Introduction to Economics}
% ============================================================

\begin{frame}{What is Economics?}

\begin{block}{Definition}
  \textbf{Economics} is the study of how individuals, businesses, governments, and societies make choices to allocate \textit{scarce} resources to satisfy their \textit{unlimited} wants and needs.
\end{block}
\vspace{0.4em}
\begin{itemize}
  \item \textbf{Scarcity} is the fundamental economic problem --- resources are limited, but wants are not
  \item Every choice involves a \textbf{trade-off}: choosing one thing means giving up something else
  \item The value of the next best alternative forgone is the \textbf{opportunity cost}
  \item Economics helps us understand \textit{how} decisions are made and \textit{what} their consequences are
\end{itemize}
\end{frame}

% ============================================================
\begin{frame}{Two Branches of Economics}
\transfade
\begin{columns}
  \begin{column}{0.48\textwidth}
    \begin{block}{Microeconomics}
      Studies the decisions of \textbf{individual} agents:
      \begin{itemize}
        \item Households and consumers
        \item Firms and producers
        \item Individual markets
      \end{itemize}
      \vspace{0.3em}
      \textit{Example: How does a rise in coffee prices affect consumer demand?}
    \end{block}
  \end{column}
  \begin{column}{0.48\textwidth}
    \begin{block}{Macroeconomics}
      Studies the economy as a \textbf{whole}:
      \begin{itemize}
        \item National output (GDP)
        \item Unemployment rates
        \item Inflation and price levels
        \item Government policy
      \end{itemize}
      \vspace{0.3em}
      \textit{Example: Why does Canada's unemployment rate rise during a recession?}
    \end{block}
  \end{column}
\end{columns}
\vspace{0.4em}
\centering \textcolor{accent}{\textbf{In ECON 101, we focus on Macroeconomics.}}
\end{frame}

% ============================================================
\begin{frame}{Why Study Macroeconomics?}

\begin{itemize}
  \item \textbf{It affects your daily life} --- interest rates on loans, job prospects, prices of goods
  \item \textbf{It shapes government decisions} --- taxes, spending, monetary policy
  \item \textbf{It explains news headlines} --- recessions, inflation, central bank rate decisions
  \item \textbf{It matters for business} --- firms plan hiring, investment, and pricing based on economic conditions
  \item \textbf{It informs your citizenship} --- understanding policy debates makes you a more informed voter and participant in society
\end{itemize}
\end{frame}

% ============================================================
\begin{frame}{Key Macroeconomic Questions}
\transfade
\begin{block}{The Big Questions We Will Answer This Semester}
  \begin{enumerate}
    \item What determines the total output and income of an economy?
    \item Why do economies grow over time --- and why do some grow faster than others?
    \item What causes unemployment, and what can be done about it?
    \item What causes inflation, and how is it controlled?
    \item How do banks and the Bank of Canada influence the economy?
    \item What role should government play through fiscal and monetary policy?
    \item How are economies around the world connected through trade and finance?
  \end{enumerate}
\end{block}
\end{frame}

% ============================================================
\begin{frame}{Positive vs.\ Normative Economics}

\begin{columns}
  \begin{column}{0.48\textwidth}
    \begin{block}{Positive Economics}
      Describes \textit{what is} --- objective statements about the world that can be tested with data.
      \vspace{0.3em}

      \textit{``A 1\% increase in the unemployment rate reduces GDP growth by approximately 2\%.''} (Okun's Law)
    \end{block}
  \end{column}
  \begin{column}{0.48\textwidth}
    \begin{block}{Normative Economics}
      Prescribes \textit{what ought to be} --- value judgments about economic outcomes or policies.
      \vspace{0.3em}

      \textit{``The government should increase the minimum wage to reduce income inequality.''}
    \end{block}
  \end{column}
\end{columns}
\vspace{0.5em}
\begin{alertblock}{Key Takeaway}
  Economists try to be \textbf{positive} in their analysis, but policy debates are often \textbf{normative}. Knowing the difference matters.
\end{alertblock}
\end{frame}

% ============================================================
\begin{frame}{Economic Models}

\begin{block}{What is an Economic Model?}
  A \textbf{model} is a simplified representation of reality used to analyze economic behaviour and make predictions.
\end{block}
\vspace{0.3em}
\begin{itemize}
  \item Models \textit{deliberately} leave out details to focus on what matters most
  \item A map is a good analogy --- it simplifies the world but is still useful for navigation
  \item We will use models involving \textbf{graphs}, \textbf{equations}, and \textbf{tables}
  \item The \textbf{ceteris paribus} assumption (``all else equal'') is central to most models --- we change one variable at a time
\end{itemize}
\vspace{0.3em}
\begin{block}{Examples We Will Use}
  Supply \& Demand model, GDP expenditure model, AD--AS model, money market model
\end{block}
\end{frame}

% ============================================================
\begin{frame}{What to Expect This Semester}
\transfade
\begin{columns}
  \begin{column}{0.48\textwidth}
    \begin{block}{To Succeed in This Course}
      \begin{itemize}
        \item Attend \textbf{all lectures} and take notes
        \item Read the textbook \textit{before} class
        \item Complete \textbf{RISE modules} by Sep 25
        \item Practice with problem sets on Moodle
        \item Visit office hours early --- don't wait until before exams
        \item Review material \textbf{regularly}, not just before tests
      \end{itemize}
    \end{block}
  \end{column}
  \begin{column}{0.48\textwidth}
    \begin{block}{Upcoming Deadlines}
      \begin{itemize}
        \item \textbf{Sep 25:} RISE Required Modules due
        \item \textbf{Sep 28:} Quiz 1 (Ch.2--3)
        \item \textbf{Oct 7:} Quiz 2 (Ch.4)
        \item \textbf{Oct 14:} Quiz 3 (Ch.5)
        \item \textbf{Oct 19:} \textbf{Midterm 1}
      \end{itemize}
    \end{block}

  \end{column}
\end{columns}
\end{frame}

% ============================================================
\begin{frame}[plain]
\transfade
  \begin{center}
    \vspace{2em}
    {\Large \textbf{Questions?}}\\[1em]
    {\normalsize karimzada@unbc.ca}\\[0.5em]
    {\normalsize Office: Room 10-4522}\\[0.5em]
    {\normalsize Office Hours: Tue 1--2:30pm $\cdot$ Wed 11:30am--1pm $\cdot$ Thu 1--2pm}\\[1.5em]
    \textcolor{accent}{\rule{0.6\textwidth}{0.4pt}}\\[1em]
    {\small See you Monday, September 14\textsuperscript{th}!}\\
    {\small Read Chapters 1 and 2 for next week.}
  \end{center}
\end{frame}

% ============================================================
\end{document}
